\documentclass[a4paper,12pt]{article}
\usepackage[T2A]{fontenc}
\usepackage[utf8]{inputenc}
\usepackage[russian]{babel}
\usepackage{amsmath,amssymb,icomma,booktabs,tabularx,array}
\usepackage[a4paper,left=20mm,right=20mm,top=23mm,bottom=20mm,headheight=25pt,headsep=7mm]{geometry}
\usepackage{xcolor,tikz,fancyhdr,enumitem,setspace,parskip,microtype,titlesec}
\usetikzlibrary{arrows.meta,positioning,shapes.geometric}
\definecolor{accent}{HTML}{008080}
\definecolor{darkaccent}{HTML}{004D4D}
\definecolor{textgray}{HTML}{333333}
\color{textgray}
\onehalfspacing
\setlength{\parindent}{0pt}
\setlength{\parskip}{5pt}
\titleformat{\section}{\Large\normalfont\color{darkaccent}}{\thesection.}{.6em}{}[\vspace{3pt}{\color{accent}\titlerule}]
\titleformat{\subsection}{\large\normalfont\color{darkaccent}}{}{0pt}{}
\titlespacing*{\section}{0pt}{14pt}{12pt}
\titlespacing*{\subsection}{0pt}{9pt}{3pt}
\pagestyle{fancy}
\fancyhf{}
\fancyhead[L]{\small Степень с натуральным показателем}
\fancyhead[R]{\small Кирилл Зиновьев}
\fancyfoot[C]{\small\thepage}
\renewcommand{\headrulewidth}{0.4pt}
\renewcommand{\headrule}{\hbox to\headwidth{\color{accent}\leaders\hrule height \headrulewidth\hfill}}
\setlist[itemize]{leftmargin=*,itemsep=1pt,topsep=2pt}
\newcommand{\idea}[1]{\textcolor{darkaccent}{\textbf{#1}}}
\newcommand{\ex}[1]{\subsection*{#1}}
\newcommand{\cb}{\(\square\)\,}
\begin{document}
\begin{titlepage}
\centering
\vspace*{27mm}
{\Large\color{darkaccent} МАТЕМАТИКА\par}
\vspace{15mm}
{\Huge Чек-лист\par}
\vspace{10mm}
{\LARGE\color{accent} Степень с натуральным показателем\\[3pt] Квадраты и кубы чисел\par}
\vspace{14mm}
{\large Вычисления и представление чисел в виде степени\par}
\vfill
{\large Кирилл Зиновьев\par}
\vspace{9mm}
{25 ноября 2026 г.\par}
\vspace{15mm}
{\small Лёвин Артём Александрович\\эксклюзивно для Кирилла Зиновьева\par}
\vspace{19mm}
\begin{tikzpicture}
\draw[accent,line width=.8pt] (0,0) .. controls (2,0.85) and (4,-.85) .. (6,0) .. controls (8,.85) and (10,-.85) .. (12,0);
\fill[darkaccent] (6,0) circle (1.1pt);
\end{tikzpicture}
\vspace*{3mm}
\end{titlepage}
\setcounter{page}{1}
\section{Что такое степень}
Степень $a^n$ с натуральным показателем $n\ge 1$ означает произведение $n$ одинаковых множителей, равных $a$.
\[
 a^n=\underbrace{a\cdot a\cdot\ldots\cdot a}_{n\text{ множителей}},\qquad a^1=a.
\]
\begin{center}
\begin{tikzpicture}[>=Stealth]
\node[font=\Huge] (power) at (0,0) {$5^3$};
\node[anchor=east,text=darkaccent] (base) at (-1.5,-.6) {основание};
\node[anchor=west,text=darkaccent] (exp) at (1.4,.8) {показатель};
\draw[->,accent] (base.east) -- (-.32,-.15);
\draw[->,accent] (exp.west) -- (.23,.35);
\end{tikzpicture}
\end{center}
Например, $5^3=5\cdot5\cdot5=125$. Показатель 3 задаёт число множителей, поэтому $5^3\ne 5\cdot3$.
\section{Знак степени и роль скобок}
Сначала определите, входит ли минус в основание. При $a>0$ и натуральном $k$:
\[
 (-a)^{2k}=a^{2k},\qquad (-a)^{2k-1}=-a^{2k-1}.
\]
\begin{center}
\renewcommand{\arraystretch}{1.3}
\begin{tabularx}{\linewidth}{@{}lXl@{}}
\toprule
Запись & Что возводится в степень & Значение \\\midrule
$-5^4$ & Только число 5 & $-625$\\
$(-5)^4$ & Число $-5$ целиком & $625$\\
$-(-5)^4$ & Число $-5$; внешний минус сохраняется & $-625$\\
$(-5)^3$ & Число $-5$, показатель нечётный & $-125$\\
\bottomrule
\end{tabularx}
\end{center}
\ex{Пример. Порядок действий}
\[
 -2^6+(-3)^4=-64+81=17.
\]
\idea{Проверяйте:} степень вычисляется прежде внешнего минуса; затем выполняются сложение и вычитание.
\clearpage
\section{Обыкновенные и десятичные дроби}
При возведении дроби в натуральную степень возводятся в степень её числитель и знаменатель:
\[
 \left(\frac{a}{b}\right)^n=\frac{a^n}{b^n},\qquad b\ne0.
\]
\ex{Пример 1. Обыкновенная дробь}
\[
 \left(-\frac12\right)^4=\frac{(-1)^4}{2^4}=\frac1{16},\qquad
 \left(-\frac23\right)^3=-\frac8{27}.
\]
В первом примере показатель чётный, во втором --- нечётный.
\ex{Пример 2. Десятичная дробь}
\[
 (-0,9)^3=\left(-\frac9{10}\right)^3=-\frac{729}{1000}=-0,729.
\]
Если у конечной десятичной дроби $d$ знаков после запятой, при возведении в степень $n$ можно сначала возвести в степень число без запятой, затем отделить справа $dn$ знаков (при необходимости дописать нули слева). Знак определите отдельно.
\section{Квадраты и кубы: опорные значения}
Квадрат числа --- его вторая степень, куб --- третья степень.
\[
 a^2=a\cdot a,\qquad a^3=a\cdot a\cdot a.
\]
\begin{center}
\renewcommand{\arraystretch}{1.25}
\begin{tabular}{@{}lrrrrrrrrrr@{}}
\toprule
$a$ &1&2&3&4&5&6&7&8&9&10\\\midrule
$a^2$ &1&4&9&16&25&36&49&64&81&100\\
$a^3$ &1&8&27&64&125&216&343&512&729&1000\\\bottomrule
\end{tabular}
\end{center}
Дополнительные квадраты для быстрого счёта:
\[
 11^2=121,\quad 12^2=144,\quad 13^2=169,\quad 14^2=196,\quad 15^2=225.
\]
\idea{Полезно запомнить:} квадрат действительного числа всегда неотрицателен; куб отрицательного числа отрицателен.
\clearpage
\section{Как представить число квадратом или кубом}
Цель --- подобрать основание и записать исходное число в виде $a^2$ или $a^3$.
\ex{Целые числа}
\[
 64=4^3,\qquad -216=(-6)^3,\qquad 169=13^2.
\]
\ex{Обыкновенные дроби}
Представьте числитель и знаменатель степенями с одинаковым показателем и объедините их:
\[
 \frac{25}{169}=\frac{5^2}{13^2}=\left(\frac5{13}\right)^2,
 \qquad -\frac1{64}=\left(-\frac14\right)^3.
\]
\ex{Смешанные числа}
Сначала переведите смешанное число в неправильную дробь, при необходимости сократите её.
\[
 1\frac{24}{25}=\frac{1\cdot25+24}{25}=\frac{49}{25}
 =\left(\frac75\right)^2.
\]
\[
 4\frac{17}{27}=\frac{4\cdot27+17}{27}=\frac{125}{27}
 =\left(\frac53\right)^3.
\]
\ex{Десятичные дроби}
Переведите число в обыкновенную дробь со знаменателем $10$, $100$, $1000$ и т.~д. Затем примените тот же метод.
\[
 0,81=\frac{81}{100}=\left(\frac9{10}\right)^2=(0,9)^2.
\]
\[
 0,0081=\frac{81}{10000}=\left(\frac9{100}\right)^2=(0,09)^2.
\]
\[
 0,008=\frac8{1000}=\left(\frac2{10}\right)^3=(0,2)^3.
\]
\idea{Важное уточнение.} Для точного представления ненулевой конечной десятичной дроби квадратом конечной десятичной дроби в её сокращённой десятичной записи число знаков после запятой должно быть чётным, а целое число без запятой --- полным квадратом. Например, $0,081$ не является квадратом конечной десятичной дроби. Это не исключает существования квадратного корня в действительных числах.
\clearpage
\section{Степени десяти}
У числа $10^n$ после единицы записано ровно $n$ нулей, если $n$ --- натуральное число.
\begin{center}
\renewcommand{\arraystretch}{1.25}
\begin{tabular}{@{}rl@{\qquad}rl@{}}
\toprule
Степень&Число&Степень&Число\\\midrule
$10^1$&$10$&$10^4$&$10\,000$\\
$10^2$&$100$&$10^5$&$100\,000$\\
$10^3$&$1000$&$10^6$&$1\,000\,000$\\\bottomrule
\end{tabular}
\end{center}
\[
 4\cdot10^3=4\cdot1000=4000,\qquad 5\cdot10^6=5\,000\,000.
\]
\section{Разложение на простые множители}
Чтобы представить целое число больше единицы в виде степени, последовательно делите его на простые делители до получения единицы. Затем подсчитайте повторяющиеся множители.
\begin{center}
\begin{tabular}{r|l@{\hspace{20mm}}r|l@{\hspace{20mm}}r|l}
125&5&32&2&216&2\\
25&5&16&2&108&2\\
5&5&8&2&54&2\\
1&&4&2&27&3\\
&&2&2&9&3\\
&&1&&3&3\\
&&&&1&
\end{tabular}
\end{center}
\[
 125=5^3,\qquad 32=2^5,\qquad
 216=2^3\cdot3^3=(2\cdot3)^3=6^3.
\]
Использовано свойство произведения:
\[
 a^nb^n=(ab)^n.
\]
\ex{Как распознать квадрат или куб по разложению}
Для целого числа $N>1$ показатели всех простых множителей должны делиться на 2, чтобы получить квадрат, или на 3, чтобы получить куб. Например:
\[
 144=2^4\cdot3^2=(2^2\cdot3)^2=12^2.
\]
\clearpage
\section*{Алгоритм: представление целого числа в виде степени}
\begin{center}
\begin{tikzpicture}[
  node distance=9mm and 11mm,
  start/.style={ellipse,draw=accent,thick,align=center,minimum width=37mm,minimum height=10mm},
  actionnode/.style={rectangle,rounded corners=2pt,draw=accent,align=center,text width=85mm,minimum height=12mm,inner sep=5pt},
  decision/.style={diamond,aspect=2.3,draw=accent,align=center,text width=53mm,inner sep=1pt},
  smallactionnode/.style={rectangle,rounded corners=2pt,draw=accent,align=center,text width=48mm,minimum height=13mm,inner sep=4pt},
  arrow/.style={-{Stealth[length=2.6mm]},thick,draw=darkaccent}]
\node[start] (s) {$N>1$, $n\ge2$ --- целые};
\node[actionnode,below=of s] (a) {Разложите $N$ на простые множители последовательным делением до 1.};
\node[actionnode,below=of a] (b) {Соберите одинаковые множители в степени: $N=p_1^{e_1}\cdot\ldots\cdot p_r^{e_r}$.};
\node[decision,below=of b] (c) {Все показатели $e_i$ делятся на $n$?};
\node[smallactionnode,below left=12mm and -11mm of c] (d) {Разделите каждый показатель на $n$.};
\node[smallactionnode,below right=12mm and -11mm of c] (e) {Целого основания для данной степени нет.};
\node[smallactionnode,below=of d] (f) {Перемножьте полученные степени: $a=p_1^{e_1/n}\cdots p_r^{e_r/n}$.};
\node[smallactionnode,below=of f] (g) {Запишите $N=a^n$. Проверьте вычислением.};
\node[start,below=of g] (h) {Готово};
\draw[arrow] (s)--(a);\draw[arrow] (a)--(b);\draw[arrow] (b)--(c);
\draw[arrow] (c.west) -| node[pos=.25,above,font=\small]{да} (d.north);
\draw[arrow] (c.east) -| node[pos=.25,above,font=\small]{нет} (e.north);
\draw[arrow] (d)--(f);\draw[arrow] (f)--(g);\draw[arrow] (g)--(h);
\draw[arrow] (e.south) |- (h.east);
\end{tikzpicture}
\end{center}
\clearpage
\section*{Тренировочный блок}
Выполняйте задания письменно. При возведении в степень определяйте знак и показывайте основные преобразования. В заданиях 9--10 используйте разложение на простые множители.
\subsection*{Базовый уровень}
\begin{enumerate}[leftmargin=*,itemsep=9pt]
\item Вычислите: $3^4$.
\item Вычислите: $-4^4+(-3)^4$.
\item Вычислите: $\displaystyle\left(-\frac25\right)^3$.
\item Вычислите: $(-0,7)^3$.
\item Вычислите: $6\cdot10^5$.
\end{enumerate}
\subsection*{Средний уровень}
\begin{enumerate}[leftmargin=*,start=6,itemsep=9pt]
\item Представьте $0,0036$ в виде квадрата положительной конечной десятичной дроби.
\item Представьте $\displaystyle 2\frac79$ в виде квадрата положительной обыкновенной дроби.
\item Представьте $-0,027$ в виде куба десятичной дроби.
\end{enumerate}
\subsection*{Повышенный уровень}
\begin{enumerate}[leftmargin=*,start=9,itemsep=9pt]
\item Представьте $243$ в виде степени числа $3$.
\item Представьте $1728$ в виде куба целого числа.
\end{enumerate}
\clearpage
\section*{Ответы}
\begin{center}
\renewcommand{\arraystretch}{1.45}
\begin{tabular}{@{}cl@{}}
\toprule
№ & Ответ\\\midrule
1 & $81$\\
2 & $-175$\\
3 & $-\dfrac8{125}$\\
4 & $-0,343$\\
5 & $600\,000$\\
6 & $(0,06)^2$\\
7 & $\left(\dfrac53\right)^2$\\
8 & $(-0,3)^3$\\
9 & $3^5$\\
10 & $12^3$\\\bottomrule
\end{tabular}
\end{center}
\section*{Итоговая памятка}
\cb Я различаю основание и показатель степени.\\
\cb Я учитываю скобки, внешний минус и чётность показателя.\\
\cb Я умею возводить в степень обыкновенные и десятичные дроби.\\
\cb Я помню основные квадраты, кубы и степени десяти.\\
\cb Я умею переводить смешанное число в неправильную дробь.\\
\cb Я умею раскладывать число на простые множители и объединять степени с одинаковыми показателями.
\end{document}
